%% Editorial completion of the representative simple-type term grammar.
\noindent\textit{Editorial convention for the preceding simple-type syntax.}
The exact source chapter fixes the base type \(o\), application, abstraction,
typed equality, and universal quantification.  The following definition makes
its subsequently used classical Boolean constants explicit; it does not select
a proof calculus or alter the source's standard/Henkin semantic distinction.

\begin{hbdefinition}{def:p02-stt-boolean-constants}{Boolean constants of the representative simple type theory}
The term grammar contains the primitive logical constants
\[
 \bottom:o,\qquad
 \neg:o\to o,\qquad
 \land,\lor,\to:o\to(o\to o),
\]
with application associating to the left.  A standard or Henkin type structure
interprets them by the classical Boolean operations on
\(D_o=\{0,1\}\): \(\bottom\) denotes \(0\), negation swaps \(0,1\), and the
three binary constants denote conjunction, disjunction, and material
implication.  Put \(\topv\defeq\neg\bottom\); for
\(P:\alpha\to o\), put \(\exists_\alpha(P)\defeq
\neg\forall_\alpha(\lambda x{:}\alpha.\neg P x)\).

These declarations are necessary because formulas are terms of type \(o\), so
untyped punctuation would not form terms.  They name object-language typed
constants, differ from the metalinguistic connectives used to describe the
grammar, and are used by the extensionality and semantic formulas that follow.
In a Henkin structure their declared denotations therefore belong to the
selected function domains at the displayed types.
\index{simple type theory!Boolean constants}
\end{hbdefinition}
